\chapter{Phụ Huynh Cấm Con Nhưng Không Cấm Mình}
\markboth{Phụ Huynh Cấm Con Nhưng Không Cấm Mình}{Phụ Huynh Cấm Con Nhưng Không Cấm Mình}

\begin{center}
	\textit{\color{bookprimary}``Đứa trẻ không học nhiều nhất từ lời dặn. Nó học từ cảnh tượng lặp đi lặp lại trước mặt nó.''}

	\textit{\color{bookprimary}``A child learns most not from warnings, but from the repeated scenes happening in front of them.''}
\end{center}

\ngancach

\section{Bài Giảng Tắt Ở Ngón Tay Người Lớn}
\begin{truyen}{Bài Giảng Tắt Ở Ngón Tay Người Lớn}{The Lecture Dies in Adult Fingers}
\chuhoa{C}{ó phụ huynh mắng con rất nặng}: ``Suốt ngày chỉ biết điện thoại. Học hành gì kiểu đó?''
\textit{(A parent scolds the child harshly: ``All you do is stay on your phone. How can you study like that?'')}

Nói xong, chính người đó lại cúi xuống trả lời tin nhắn ngay trong lúc con còn đang ngồi trước mặt.
\textit{(Right after saying that, the same parent bends down to reply to a message while the child is still sitting there.)}

\adultpair{Đó là một trong những cảnh làm giáo dục mất điện nhanh nhất. Người lớn không thể đòi một đứa trẻ tôn trọng điều mình đang phá ngay trước mắt nó. Trẻ con có thể không cãi thành lời, nhưng trong đầu nó sẽ ghi rất rõ: điện thoại của người lớn là việc quan trọng; điện thoại của mình thì bị coi là tội.}{``That is one of the fastest ways to drain all power from education. Adults cannot demand that a child respect something they themselves are breaking right in front of that child. Children may not argue aloud, but in their minds they record it clearly: the adult's phone is important business; mine is treated as a sin.''}

Người lớn thường bảo con thiếu kỷ luật.
\textit{(Adults often say children lack discipline.)}

Nhưng rất nhiều gia đình đang thiếu thứ còn quan trọng hơn.
\textit{(But many families are lacking something even more important.)}

Họ thiếu sự nhất quán.
\textit{(They lack consistency.)}
\end{truyen}

\section{Giờ Quy Định Chỉ Áp Cho Con}
\begin{truyen}{Giờ Quy Định Chỉ Áp Cho Con}{Rules That Apply Only to the Child}
\chuhoa{C}{ó nhà đặt giờ cất điện thoại cho con rất rõ, nhưng người lớn thì muốn dùng lúc nào cũng được.}
\textit{(Some homes set clear phone hours for the child, while adults reserve the right to use theirs whenever they want.)}

Đứa trẻ nhìn thấy một điều đơn giản.
\textit{(The child sees one simple thing.)}

Luật ở đây không dựa trên điều đúng, mà dựa trên ai có quyền hơn.
\textit{(Rules here are not based on what is right, but on who has more power.)}

\adultpair{Một quy tắc bất công có thể vẫn được tuân bằng sợ hãi, nhưng rất khó được tôn trọng bằng lòng tin. Giáo dục bền không sống bằng ép buộc mãi. Nó sống bằng sự công bằng đủ để lời nói còn giữ được phẩm giá.}{``An unfair rule may still be obeyed out of fear, but it is rarely respected in trust. Lasting education does not survive on force alone. It survives on fairness strong enough to preserve the dignity of words.''}
\end{truyen}

\section{Bữa Cơm Cấm Máy Nhưng Bố Mẹ Ngoại Lệ}
\begin{truyen}{Bữa Cơm Cấm Máy Nhưng Bố Mẹ Ngoại Lệ}{A No-Phone Meal Except for the Parents}
\chuhoa{C}{ó gia đình yêu cầu con không mang điện thoại lên bàn ăn, còn cha mẹ thì vẫn để máy cạnh tay vì lỡ có việc.}
\textit{(A family tells the child not to bring a phone to the table, while the parents keep theirs nearby just in case.)}

Chính chữ ``ngoại lệ'' đó làm bài học xẹp xuống.
\textit{(That word ``exception'' deflates the lesson.)}

\adultpair{Nếu bữa cơm là nơi gìn giữ hiện diện, thì người lớn phải là người giữ điều đó đầu tiên. Không thể vừa dạy con trân trọng cuộc gặp vừa cầm sẵn cửa thoát khỏi cuộc gặp trên tay.}{``If mealtime is meant to protect presence, adults must protect it first. You cannot teach a child to honor a meeting while holding an exit from that meeting in your hand.''}
\end{truyen}

\section{Mượn Công Việc Để Hợp Thức Hóa}
\begin{truyen}{Mượn Công Việc Để Hợp Thức Hóa}{Using Work as a Moral Excuse}
\chuhoa{C}{ó nhiều người lớn nói rất nhanh: ``Bố mẹ khác, bố mẹ còn công việc.''}
\textit{(Many adults quickly say, ``Parents are different. We have work.'' )}

Có khi đúng thật.
\textit{(That is sometimes true.)}

Nhưng cũng có khi chữ ``công việc'' chỉ là tấm bình phong che cho một thói quen đã thành nghiện nhẹ.
\textit{(But sometimes ``work'' becomes a curtain hiding a mild addiction.)}

\adultpair{Không phải cứ liên quan công việc là mặc nhiên vô tội. Người lớn tử tế phải đủ tỉnh để phân biệt đâu là việc thật sự cần xử lý, đâu chỉ là cớ đẹp để tiếp tục bị kéo đi.}{``Not everything labeled work is automatically innocent. Decent adults must be lucid enough to distinguish what truly needs handling from what is merely a respectable excuse to keep being dragged along.''}
\end{truyen}

\section{Bảo Con Đọc Sách Nhưng Mình Lướt Liên Tục}
\begin{truyen}{Bảo Con Đọc Sách Nhưng Mình Lướt Liên Tục}{Telling the Child to Read While Scrolling Nonstop}
\chuhoa{C}{ó phụ huynh giục con ra bàn học, mở sách ra, rồi chính mình ngồi cạnh lướt video hết cái này tới cái khác.}
\textit{(A parent tells the child to study and open the book, then sits nearby scrolling through videos one after another.)}

Không khí đó rất khó sinh ra sự học nghiêm túc.
\textit{(That atmosphere hardly generates serious learning.)}

\adultpair{Trẻ con không chỉ học bằng mệnh lệnh. Nó học bằng bầu khí. Một căn nhà mà người lớn luôn phát ra nhịp phân tán thì khó đòi trẻ giữ được nhịp sâu.}{``Children do not learn only through commands. They learn through atmosphere. In a home where adults constantly radiate distraction, it is hard to demand depth from the young.''}
\end{truyen}

\section{Phạt Con Vì Mất Tập Trung Nhưng Nhà Thì Ồn Màn Hình}
\begin{truyen}{Phạt Con Vì Mất Tập Trung Nhưng Nhà Thì Ồn Màn Hình}{Punishing Distraction in a House Full of Screen Noise}
\chuhoa{C}{ó đứa trẻ bị trách vì học trước ti vi, trong khi chính phòng khách luôn mở tiếng video, tiếng tin nhắn, tiếng clip ngắn.}
\textit{(A child is scolded for studying poorly, while the house itself is filled with television, message sounds, and short clips.)}

Muốn một đứa trẻ tập trung, môi trường cũng phải đỡ phá nó.
\textit{(If a child is expected to focus, the environment must stop sabotaging that focus.)}

\adultpair{Kỷ luật không chỉ là nhắc nhở cá nhân. Còn là cách cả nhà bố trí đời sống để điều đúng có cơ hội xảy ra. Nếu nhà luôn ồn bằng màn hình, đừng ngạc nhiên khi sự chú ý của trẻ không đứng nổi.}{``Discipline is not only personal reminders. It is also the way a whole household arranges life so that the right thing has a chance to happen. If the home is constantly noisy with screens, do not be surprised when a child's attention cannot stand upright.''}
\end{truyen}

\section{Xin Con Nghe Lời Nhưng Không Sống Làm Gương}
\begin{truyen}{Xin Con Nghe Lời Nhưng Không Sống Làm Gương}{Asking for Obedience Without Living as an Example}
\chuhoa{C}{ó người lớn rất muốn con nghe mình, nhưng lại không chịu để đời sống mình bị soi ngược bởi chính lời dạy đó.}
\textit{(Some adults want children to listen, yet refuse to let their own life be examined by the advice they give.)}

Điều này làm giáo dục thành một chiều áp xuống.
\textit{(That turns education into a one-way pressure from above.)}

\adultpair{Uy quyền thật không sợ làm gương. Chỉ uy quyền rỗng mới luôn cần một vùng miễn trừ cho bản thân.}{``Real authority is not afraid of exemplarity. Only hollow authority constantly demands an exempt zone for itself.''}
\end{truyen}

\section{Một Nhà Muốn Đổi Phải Đổi Từ Người Lớn}
\begin{truyen}{Một Nhà Muốn Đổi Phải Đổi Từ Người Lớn}{If a Home Wants Change, Adults Must Start}
\chuhoa{C}{ó gia đình quyết định trước khi siết con thì cha mẹ tự cất điện thoại từ 8 giờ tối trong một tháng.}
\textit{(A family decides that before tightening rules for the child, the parents themselves will put phones away from 8 p.m. for one month.)}

Không khí nhà đổi đi thấy rõ.
\textit{(The atmosphere of the home changes noticeably.)}

Đứa trẻ ít chống hơn, vì lần đầu nó thấy quy tắc không chỉ là thứ ném xuống đầu mình.
\textit{(The child resists less, because for the first time the rule is not something thrown only at them.)}

\adultpair{Muốn sửa một nếp sống gia đình, người lớn phải chấp nhận là người trả giá đầu tiên. Không có con đường tử tế nào khác.}{``If adults want to repair a household habit, they must accept being the first to pay the price. There is no more honorable path.''}
\end{truyen}

\begin{baihoc}
Muốn dạy con rời điện thoại, trước hết người lớn phải bớt sống như thể mình không thể rời nó.
\textit{(If adults want to teach children to step away from the phone, they must first stop living as if they themselves cannot step away from it.)}

Một quy tắc chỉ có sức nặng khi người nói cũng chịu đứng trong nó.
\textit{(A rule only carries weight when the one who speaks it is willing to stand inside it too.)}
\end{baihoc}

\begin{tuhocnhanh}
1. Trong nhà em, quy tắc dùng điện thoại có áp cho mọi người như nhau không? \textit{(In your home, do phone rules apply equally to everyone?)}

2. Người lớn quanh em đang dạy bằng lời nhiều hơn hay bằng nếp sống nhiều hơn? \textit{(Are the adults around you teaching more by words or by way of living?)}

3. Nếu em là cha mẹ, em có chịu được việc chính con mình soi lại thói quen của mình không? \textit{(If you were a parent, could you bear your own child examining your habits?)}
\end{tuhocnhanh}

\begin{tongketchuong}
Trẻ con không cần người lớn hoàn hảo.
\textit{(Children do not need perfect adults.)}

Nhưng chúng rất cần người lớn đừng sống giả ngay trong lúc dạy dỗ.
\textit{(But they desperately need adults not to live hypocritically while teaching.)}
\end{tongketchuong}

\ngancach